\documentclass[a4paper]{article}

\usepackage[fontset=fandol]{ctex}
\usepackage{amsmath, amssymb}
\usepackage{graphicx}
\usepackage[margin=2.5cm]{geometry}
\usepackage[most]{tcolorbox}
\usepackage{etoolbox}
\usepackage{listings}
\usepackage{hyperref}
\usepackage{booktabs}
\usepackage{subcaption}
\usepackage{float}
\usepackage{tikz}
\IfFileExists{pgfplots.sty}{
    \usepackage{pgfplots}
    \pgfplotsset{compat=1.18}
}{}

\newtcolorbox{knowledgebox}[1]{
    enhanced, colback=blue!5!white, colframe=blue!75!black, colbacktitle=blue!75!black,
    coltitle=white, fonttitle=\bfseries, title=#1, attach boxed title to top left={yshift=-2mm, xshift=2mm},
    boxrule=1pt, sharp corners
}

\newtcolorbox{importantbox}[1]{
    enhanced, colback=yellow!10!white, colframe=yellow!80!black, colbacktitle=yellow!80!black,
    coltitle=black, fonttitle=\bfseries, title=#1, sharp corners
}

\newtcolorbox{warningbox}[1]{
    enhanced, colback=red!5!white, colframe=red!75!black, colbacktitle=red!75!black,
    coltitle=white, fonttitle=\bfseries, title=#1, sharp corners
}

\lstset{
    language=Python,
    basicstyle=\ttfamily\small,
    keywordstyle=\color{blue},
    stringstyle=\color{red!60!black},
    commentstyle=\color{green!60!black},
    breaklines=true,
    frame=single,
    numbers=left,
    numberstyle=\tiny\color{gray},
    captionpos=b,
    extendedchars=false
}

\newcommand{\notetitle}{CS294-277 课程笔记：Lecture 11B 灵巧操作案例（下）}
\newcommand{\noteauthors}{Codex}
\newcommand{\notedate}{2026-04-16}
\newcommand{\videochannel}{CS294-277 / UC Berkeley}
\newcommand{\videopublishdate}{2024-11-25}
\newcommand{\videoduration}{未在本地元数据中提供}
\newcommand{\videourl}{https://www.youtube.com/watch?v=mkI7WGgj5xA}
\newcommand{\videocoverpath}{cover.jpg}

\begin{document}

\begin{titlepage}
\centering
{\Large 课程笔记\par}
\vspace{1.2cm}
{\huge\bfseries \notetitle\par}
\vspace{0.8cm}
{\large \noteauthors\par}
\vspace{0.3cm}
{\large \notedate\par}
\vspace{1.2cm}

\ifdefempty{\videocoverpath}{
{\small 请在生成笔记时填入视频封面图的本地路径。\par}
}{
\includegraphics[width=0.82\textwidth,height=0.45\textheight,keepaspectratio]{\videocoverpath}\par
}

\vfill
\begin{tcolorbox}[width=0.9\textwidth, colback=black!2!white, colframe=black!60, sharp corners]
\textbf{视频作者/频道}：\videochannel\par
\textbf{发布日期}：\videopublishdate\par
\textbf{视频时长}：\videoduration\par
\textbf{视频链接}：\href{\videourl}{\nolinkurl{\videourl}}
\end{tcolorbox}
\end{titlepage}

\tableofcontents
\newpage

\begin{figure}[H]
\centering
\includegraphics[width=0.78\textwidth]{overview_cover.jpg}
\caption{课程官方材料中的代表性页面，用于帮助读者在进入本讲正文前建立整体上下文。}
\end{figure}
\section{本讲定位与核心矛盾}

这节课延续 11A 的灵巧操作主题，但关注点从“监督学习怎么做”切换到了“当任务拉长到多阶段、多对象、多场景时，系统如何仍然能稳定工作”。讲者选取的主案例是 Dalal et al. (2024) 的 ManipGen，它试图让一个单一 generalist agent 在真实场景中完成 zero-shot long-horizon manipulation。

\subsection{为什么长时序操控比单步操控更难}

长时序任务的难点不只是动作序列更长，而是任务会跨越多个层次：
\begin{itemize}
  \item 先要理解任务目标；
  \item 再要决定子任务顺序；
  \item 然后要在局部接触中完成稳定控制；
  \item 最后还要在失败时恢复并继续前进。
\end{itemize}

课程摘录指出，若想靠纯数据规模解决这类问题，数据集往往需要增长到 100K--1M trajectories 的量级，这对人工采集几乎不可持续。另一方面，sim2real（simulation-to-reality）也不是自然成立的，因为接触动力学、场景组合和长时序行为都很难精确模拟。

\subsection{本讲的中心假设}

\begin{importantbox}{ManipGen 的关键判断}
许多操作任务可以分解为“接近目标”和“在目标附近发生接触”两个阶段。只要把第二阶段限制在局部坐标系和局部观测里，许多看似复杂的问题就会突然变得更可学。
\end{importantbox}

\subsection{本章小结}

本讲的主线不是“更多端到端学习”，而是“更合理的分层与局部化”。ManipGen 的整个系统设计都围绕这个原则展开。

\section{ManipGen 的局部性假设}

\subsection{从全局任务到局部交互}

很多家庭和工业操作任务都可以拆成两步。比如开门：
\begin{enumerate}
  \item 先把手移动到门把附近的预接触位姿；
  \item 再在门把附近执行抓握、旋转、拉动等接触丰富的动作。
\end{enumerate}

对第二步来说，机器人并不需要知道整个房间的几何细节，只要知道门把和手的相对关系、附近局部视觉信息以及动作限制即可。这就是 locality of interaction（交互局部性）。

\subsection{局部策略的定义}

讲义给出的核心形式是局部策略（local policy）：
\[
\pi_{\mathrm{loc}}(a^{t}_{\mathrm{loc}} \mid O^{t}_{\mathrm{loc}})
\]
它只看局部观测 \(O^{t}_{\mathrm{loc}}\)，并输出局部动作 \(a^{t}_{\mathrm{loc}}\)。如果初始化时末端执行器已经足够接近目标，那么局部策略就可以忽略绝对位姿，把注意力放在相对几何和接触上。

\subsection{局部策略的三个不变性}

课程摘录强调，局部策略对三类变化更鲁棒：
\begin{itemize}
  \item \textbf{Absolute pose}：只关心相对位置，不依赖全局位姿；
  \item \textbf{Skill order}：任务链可以由外部 motion planning 或高层策略组织；
  \item \textbf{Scene configuration}：只观察目标附近的局部区域，不必建模整场景。
\end{itemize}

\subsection{本章小结}

ManipGen 的重要贡献，不是“把所有事情都学进一个大模型”，而是把长时序操控拆成可重复调用的局部模块。

\section{从 RL 学局部技能}

\subsection{局部 skill library}

ManipGen 的第一步是先通过强化学习（RL）在模拟中学到一组局部技能库。讲义列出的五个技能是：
\begin{itemize}
  \item pick
  \item place
  \item grasp handle
  \item open
  \item close
\end{itemize}

这些技能之所以可以单独学习，是因为每个技能都能被局部状态、局部目标和局部奖励描述。

\subsection{奖励设计与 PPO}

课程摘录给出一个复合奖励形式：
\[
r = c_1 r_{\mathrm{ee}} + c_2 r_{\mathrm{obj}} + c_3 r_{\mathrm{ee,obj}} + c_4 r_{\mathrm{action}} + c_5 r_{\mathrm{succ}}
\]
其中 \(r_{\mathrm{ee}}\)、\(r_{\mathrm{obj}}\) 和 \(r_{\mathrm{succ}}\) 分别强调末端执行器、目标对象以及成功性；\(r_{\mathrm{ee,obj}}\) 和 \(r_{\mathrm{action}}\) 则用于约束运动与便于 sim2real。

这些局部策略使用 privileged observation 和 proprioceptive state 训练，并采用 PPO（Proximal Policy Optimization, PPO）作为优化器。这样做的好处是，RL 在简单局部技能上可以提供密集反馈，而且不需要大量人工演示。

\subsection{为什么这一步适合用 RL}

局部技能的关键不在于全局任务多么复杂，而在于局部接触是否可以被稳定学习。只要局部状态足够简洁，RL 的探索和优化就比从人类示范里直接学一个超长行为序列更自然。

\subsection{本章小结}

ManipGen 先把“难的长任务”拆成“容易学的局部技能”，再用 RL 把这些技能做扎实。这是整个方法能成立的第一层基础。

\section{从专家到通用策略的蒸馏}

\subsection{多任务 DAgger}

在学到局部专家后，ManipGen 进一步用在线模仿学习（online imitation learning）和多任务 DAgger（Dataset Aggregation, DAgger）把这些单任务局部专家蒸馏成更一般的 visuomotor policy。

这里的逻辑很重要：不是直接拿一个通用策略从头学，而是先让局部专家在各自的窄域内做到可靠，再把它们的行为知识迁移到更一般的策略中。

\subsection{观测设计：为什么只用局部信息}

蒸馏阶段的观测设计并不是“越多越好”，而是“与局部目标最相关的信息越够用越好”。课程摘录提到，ManipGen 主要使用：
\begin{itemize}
  \item local wrist camera information；
  \item current frame image；
  \item first frame depth image；
  \item first frame semantic image。
\end{itemize}

其中 Grounded SAM（Grounded Segment Anything Model）被用来提取语义信息，从而帮助定位目标对象。这样一来，系统并不需要依赖过重的端到端视觉编码，而是可以把一部分视觉理解交给成熟的计算机视觉模块。

\subsection{网络结构}

蒸馏后端使用的是 ResNet18 + MLP 的组合。这个选择本身也说明了一种工程原则：当高层任务已经通过结构化分解被简化时，末端策略并不一定非要使用最复杂的网络。

\subsection{本章小结}

蒸馏阶段的本质，是把“局部专家能做”的知识转成“通用策略会做”的参数化表达，同时保留局部观测和任务分解的结构优势。

\section{部署：VLM、运动规划与局部策略}

\subsection{Text-conditioned manipulation}

ManipGen 的最终部署不是一个单模型直接端到端出动作，而是三段式的 hierarchical decomposition：
\begin{enumerate}
  \item 用 VLM（vision-language model）根据场景和文本提示做高层规划；
  \item 用 neural motion planning 把机器人带到局部可操作位姿；
  \item 用 local policies 完成局部接触动作；
\end{enumerate}
然后把这三段循环起来，直到完成整段长时序任务。

\subsection{为什么语言在这里有用}

讲义明确指出，language model 可能并不擅长直接输出低层动作，但它在高层规划上很有用。原因是语言更适合表达目标、子步骤和操作约束，而不是毫秒级的控制细节。

\subsection{Zero-shot performance}

课程摘录给出的结果显示，ManipGen 在多个真实任务上取得了 76\% 的 zero-shot success rate，平均完成 4.28 / 4.8 个阶段。相比之下，OpenVLA、SayCan、LLMTrajGen 和 VoxPoser 在这些任务上的表现更弱。

这类结果真正说明的是：当任务足够长、足够组合化时，层级式分解和局部技能库往往比单一大模型更可靠。

\subsection{本章小结}

部署阶段证明了一个现实判断：高层规划、运动规划和局部控制各做各的，反而比试图让一个模型一次性处理所有层级更稳。

\section{类似案例：AutoMate}

\subsection{为什么它和 ManipGen 是同一类思路}

讲义最后提到 AutoMate2，说明类似的“先学局部专家，再蒸馏和组合”的策略也能用于 assembly 和 insertion 一类任务。这里的关键仍然是：RL 更适合窄空间、清晰 privileged information 的局部问题，而 imitation learning 更适合把这些局部能力推广成更一般的策略。

\subsection{几个值得记住的想法}

课程摘录提到三点很有意思：
\begin{itemize}
  \item dis-assembly 可以反向用来得到 assembly actions；
  \item human demonstrations 可以作为 RL 的初始化；
  \item dynamic time wrapping（动态时间规整）可以帮助对齐人类与机器人轨迹。
\end{itemize}

这些想法都在说明一个共同点：当任务复杂度上升时，最有效的办法通常不是直接暴力拟合，而是把问题拆成可对齐、可逆、可蒸馏的子问题。

\subsection{本章小结}

AutoMate 不是在讲一个完全不同的故事，而是在用另一个任务家族证明同一个观点：层级式、模块化的学习路线更适合真实机器人长时序操作。

\section{总结与延伸}

这一讲把灵巧操作的讨论从“单步学习”推进到了“长时序系统设计”。ManipGen 的核心启发可以概括为三句话：
\begin{itemize}
  \item 先把任务拆成局部接触和高层规划；
  \item 先用 RL 学局部专家，再用 DAgger 蒸馏到通用策略；
  \item 部署时再把 VLM、运动规划和局部策略串起来。
\end{itemize}

如果把 11A 和 11B 合起来看，这两讲其实共同建立了一个更成熟的机器人学习观：真正可扩展的系统，很少是纯端到端的，而更像是“局部技能 + 高层组织 + 适当感知”的组合体。

\subsection{本章小结}

长时序操控的关键，不是让一个大模型替代所有模块，而是让不同模块各自承担它最擅长的层次。

\end{document}
